%% front/map.tex -- the concept map. Names parts and chapters; never
%% entry ids (rule R22: ids live beside headings and in the registry,
%% the reader should not have to memorise them).

\chapter*{The Map}
\hbtopicmark{THE MAP}

This book is one subject: how people move other people, how it is done
well, how it is spotted, and what it costs the mover. The seven parts run
in the order a campaign runs --- learn the ground, read the person, take a
position, make the approach, apply or resist force, and pay the bill.

\section*{The chain}

\begin{itemize}
\item \textbf{Part I, The Terrain} --- what manipulation is before anyone
  does it: the definitions that matter, the equipment a practitioner
  actually needs, the character types walking around, and the load-bearing
  points every person can be moved from.
\item \textbf{Part II, Reading} --- the intake side: motives behind words,
  tells that would survive a camera, listening as an instrument, and
  judging who you are actually dealing with before you engage.
\item \textbf{Part III, Positioning} --- the state you arrange before the
  approach: reputation, attention and scarcity, dependence, patrons and
  alliances, concealment.
\item \textbf{Part IV, Persuasion} --- moving people without force:
  framing and controlled choice, the appeals, demonstration over argument,
  delivery, and appearance.
\item \textbf{Part V, Pressure} --- moving people with force: guilt,
  fear and obligation; coercion; open conflict; and seduction, which is
  pressure wearing soft gloves.
\item \textbf{Part VI, Defence} --- the same machinery from the far side:
  detecting it live, refusing without winning, the money and contract
  traps where real losses land, and getting out clean.
\item \textbf{Part VII, Cost} --- what the method pays out and what it
  charges: the bill for using people, and the self-mastery that is the
  precondition for surviving your own tricks.
\end{itemize}

\section*{How the parts feed each other}

Reading (Part~II) feeds everything --- every tell in Parts~IV and~V is
also a detector for Part~VI. Positioning (Part~III) is the substrate the
approaches stand on: a move made from reputation lands differently from
the same move made without one. Defence (Part~VI) mirrors Parts~IV
and~V move for move --- nearly every entry there has a twin here --- and
Part~VII prices all of it, which is why its first chapter is quoted
throughout the book rather than saved for the end.

\section*{Two ways through}

\textbf{If you are reading to defend:} start at Part~VI, then read
Part~II, then Part~VII. Detection without reading skill produces paranoid
pattern-matching; detection without the cost chapter produces a person
who defends against everything and falls for the thing they wanted to
hear.

\textbf{If you are reading the craft:} read in order. The parts assume
the vocabulary of the parts before them, and the entry format assumes you
have read \emph{How to Read an Entry}.

\section*{Finding things}

The table of contents lists every entry by title under its chapter, and
each heading carries its registry number for cross-referencing the
appendix and the source files. Cross-references inside entries are
written out --- you will never be asked to remember what a number means.
